\ssscp{\tsc{Seko-Badaic}}{C-09-05-01}
The following languages are all Austronesian,
belonging to the \tsc{Seko-Badaic} subgroup,
which may be included within the \tsc{South Sulawesi} subgroup
(\citealp[113]{Mead2003a}, \citealp[301]{Zobel2020}).
The classification of \ili{Rampi} is particularly uncertain.
\citet{Zobel2020} suggests \ili{Rampi} is an isolate coordinate with
\tsc{Celebic} and \tsc{South Sulawesi}.
Similarly, \ili{Lemolang} may not belong in this group.

\noindent{\wg\st{3pt}\begin{tabularx}{\linewidth}
{lllll>{\raggedright\arraybackslash}X}
\lsptoprule
%Language	&	ISO	&	Term	&	Gloss	&	Etymon	&	Source	\\	\midrule
%\endfirsthead
Language	&	ISO	&	Term	&	Gloss	&	Etymon	&	Source	\\	\midrule
%\endhead
\ili{Rampi}	&	lje	&	‹kühe›	&	cuscus	&	**kusay	&	Michael Martens (DM)	\\
\ili{Napu}	&	npy	&	kuhe;	&	bear cuscus;	&	**kusay	&	Michael Martens (DM)	\\	\hhline{~}
	&		&	taŋali	&	dwarf cuscus	&	**taŋali	&	\hp{\,}	\\
Behoa	&	bep	&	kuhe	&	bear cuscus	&	**kusay	&	Michael Martens (DM)	\\	\hhline{~}
	&		&	taŋali	&	dwarf cuscus	&	**taŋali	&	\hp{\,}	\\
\ili{Bada}	&	bhz	&	kuhe	&	cuscus	&	**kusay	&	Michael Martens (DM)	\\
\ili{Lemolang}	&	ley	&	kuse	&	bear cuscus	&	**kusay	&	William McConvell pers. comm. 9 June 2021	\\
\ili{Seko Padang}	&	skx	&	kuse;	&	bear cuscus;	&	**kusay	&	Michael Martens (DM)	\\	\hhline{~}
	&		&	tunnaliʔ	&	dwarf cuscus	&	(**taŋali)	&	\hp{\,}	\\
\lspbottomrule
\end{tabularx}}

%\noindent{\wg\st{2pt}\begin{xltabular}{\linewidth}
%{lllll>{\raggedright\arraybackslash}X}
%\lsptoprule
%
%\lspbottomrule
%\end{xltabular}}